\chapter{Network Traffic Analysis: Foundations}
\label{ch:foundations}

This chapter describes what a network intrusion detector observes and how its results are
judged. It covers the levels at which traffic is observed, how traffic is captured and exported,
the features built from it, the classical detection methods and the standards around them. It
ends with the flaws that recur in the evaluation of detectors.

% ================================================================
% 2.1 Network traffic: what is observed and at which layer
% ================================================================
\section{What a monitor observes}
\label{sec:network-observation-levels}

Every detection approach discussed in this thesis, whether signature-based,
purely data-driven or neurosymbolic, works on a representation of network
traffic and never on the traffic itself. That representation is chosen at a
specific level of granularity and captured at a specific point in the
network. The choice is not a neutral implementation detail. Before any
detection logic is designed, it fixes which attacks can be observed at all,
how much state the monitor has to keep, and how much the raw packet stream
is reduced. These three properties come back throughout the thesis, from the
limits of classical detection in \cref{ch:introduction} to the feature sets
used in the benchmark. This section introduces the levels of observation and
then places them on the TCP/IP stack, as seen by a passive monitor.

\subsection{Levels of observation}
\label{sec:levels-packet-flow-session-host}

Four levels of observation recur in the monitoring and intrusion-detection
literature. They are not mutually exclusive. A flow record is computed from
packets, a session groups flows or connections, and a host-level view
aggregates many sessions of the same endpoint. What changes from one level
to the next is the unit of analysis handed to the detector, and therefore
what it can and cannot see.

\paragraph{Packet level.}
The finest unit is the individual packet: its link-, network- and
transport-layer headers and, when available and unencrypted, its payload.
Signature-based detection compares units of activity such as a packet to a
list of signatures using string comparison \citep{scarfone2007nist80094}.
Snort works at this level \citep{roesch1999snort}. The price is cost.
Inspecting the content of every packet cannot easily be done at high speed
\citep{sperotto2010flowbased}.

\paragraph{Flow level.}
A \emph{flow} groups packets that share common properties into one record
and discards the packets once they have been counted. This thesis uses the
definition of the IETF \gls{ipfix} working group: a flow is ``a set of
packets or frames passing an Observation Point in the network during a
certain time interval'' \citep{claise2013rfc7011}. Membership is decided by
the \emph{Flow Key}, the fields that define the properties common to all
packets of the flow. A common choice is the 5-tuple of source and
destination address, source and destination port, and transport protocol.
The \emph{Metering Process} turns observed packets into flow records and
exports a \emph{Flow Record} for each flow. A flow is expired and exported
when it has been idle for longer than an inactivity timeout, when it reaches
a maximum lifetime, or when a TCP FIN or RST is seen
\citep{sperotto2010flowbased}. NetFlow v9 also exports long-lived flows
periodically \citep{claise2004rfc3954}.

Two properties of flows matter for the scope of this thesis. First, flows do
not carry any payload \citep{sperotto2010flowbased}. They record that two
endpoints communicated and how much, not what they exchanged. This is the
boundary set in \cref{sec:scope-limitations}: the thesis does not inspect
payload. Only on NSL-KDD do the detectors also see content features, which
come with that dataset. Second,
flow analysis reduces the work of the monitor considerably.
\citet{sperotto2010flowbased} start from the observation that packet
inspection ``cannot easily be performed at high-speeds''. In flow-based
detection ``the flow of data through the network is analyzed, instead of the
contents of each individual packet''. For the network of the University of
Twente they calculate that the packets exported by NetFlow are on average
0.1\% of the packets on the network, and that the overhead in bytes is on
average 0.2\%. Flow analysis therefore scales to fast links, at the cost of
the payload-level detail of the packet level. The same trade-off,
information lost in exchange for tractability, motivates the search for
detectors that compensate with domain knowledge, which is the subject of
\cref{ch:learning-paradigms}.

A flow should not be confused with a transport-layer connection. A UDP
exchange between two address and port pairs generates a flow just as a TCP
session does. The IPFIX definition also allows a flow to represent ``zero or
more packets'', so a single packet is enough to create one
\citep{claise2013rfc7011,sperotto2010flowbased}.

\paragraph{Session level.}
A flow is defined only by shared header properties within a metering
process. This thesis uses the word \emph{session}, or \emph{connection}, for
a bounded communication episode with an explicit beginning and end. The
classical session-level data come from the DARPA and KDD datasets, on which
much of the early \gls{ml}-based detection work and the NSL-KDD benchmark
used later in this thesis are built. \citet{tavallaee2009kddcup99} describe
the data as single connection vectors, each with 41 features and a size of
about 100 bytes, labelled as normal or as a specific attack. The attacks
fall into four classes: denial of service (DoS), probing, Remote-to-Local
(R2L) and User-to-Root (U2R). The features fall into three categories. Basic
features are the attributes that can be extracted from a TCP/IP connection.
Content features come from
inspecting the data portion of the packets. Traffic features are computed
over a window of related connections. Part of a connection record therefore
comes from content, which a bare IPFIX record does not carry. The session
level is the one at which the NSL-KDD feature set is defined, and
\cref{sec:feature-engineering-traffic} goes through its categories.

\paragraph{Host level.}
The coarsest level aggregates traffic around an endpoint, usually a
destination host and sometimes a destination service, over a window fixed
either by time or by a number of preceding records. It exists because some
attacks are invisible in any single connection and show up only in how many
connections a host is involved in. A scan of many destination hosts, for
example, produces no anomalous connection, only an unusually large number of
short connections converging on or leaving one endpoint. The KDD feature set
uses two windows for this purpose. The time-based traffic features look at
the last two seconds. Some slow probing attacks scan at much longer
intervals, for example one connection per minute, and leave no pattern in
such a window. The connection-based traffic features therefore look at the
last 100 connections \citep{tavallaee2009kddcup99}. Host-level aggregation is a second-order construction: it presupposes
that flows or connections have been extracted and computes statistics across
them. The same idea, generalised from hand-picked aggregates to explicit
domain rules, motivates the formalisation of host- and flow-level conditions
of each attack family as candidate logical axioms in
\cref{ch:attacks-ddos}.

\subsection{The TCP/IP stack seen by a monitor}
\label{sec:tcpip-stack-monitor-view}

The four levels of observation follow the layering of the protocol suite
that produces the traffic. The layering shows which fields are
available for a Flow Key, a connection record or a host aggregate, and where
payload begins. This thesis uses the four-layer description of the Internet
protocol suite in the \gls{ietf} host-requirements standard. The
\emph{application layer} is ``the top layer of the Internet protocol suite''
and hosts user protocols such as Telnet, FTP or SMTP. The \emph{transport
layer} ``provides end-to-end communication services for applications'',
either as connection-oriented, reliable \gls{tcp} or as connectionless
\gls{udp}. In the \emph{internet layer}, every Internet transport protocol
uses \gls{ip} ``to carry data from source host to destination host'', which
is ``a connectionless or datagram internetwork service, providing no
end-to-end delivery guarantees''. The \emph{link layer} is the protocol a
host ``must implement'' in order ``to communicate on its directly-connected
network'' \citep{braden1989rfc1122}.

Two aspects of this layering decide what a passive monitor at an
Observation Point can see.

The first is \emph{encapsulation}. Each layer's data unit travels unchanged
inside the data unit of the layer below. A link-layer \emph{frame} is a
link-layer header followed by a packet, an \gls{ip} \emph{datagram} is an
\gls{ip} header followed by transport-layer data, and a \gls{tcp}
\emph{segment} is a \gls{tcp} header followed by application data
\citep{braden1989rfc1122}. A monitor therefore captures nested headers, not
independent layers. The Flow Key fields come from two different headers: the
addresses and the protocol number from the \gls{ip} header, the ports from
the \gls{tcp} or \gls{udp} header. Everything after the transport-layer
header is application payload.

The second is where \emph{connection semantics} live, which settles the
difference between flow and connection. The internet layer is connectionless.
\gls{ip} moves each datagram independently, so nothing at that layer
separates an ongoing exchange from a single unrelated datagram
\citep{braden1989rfc1122}. The handshake, the sequence numbers and the
explicit close with FIN or RST, which a Metering Process uses as one of its
expiry conditions, are built by \gls{tcp} at the transport layer on top of a
network layer that knows nothing about them. \gls{udp} builds none of it
\citep{braden1989rfc1122}. A flow can therefore be computed from internet-
and transport-layer headers alone, whether or not the transport protocol
keeps connection state. A connection exists only where that state exists,
which is why the notion is meaningful for \gls{tcp} and not, in the same
sense, for \gls{udp}.

The layering also marks what the detectors in this thesis are given. Every
field used at the packet, flow, session and host levels comes from link,
internet or transport headers. Only the content features reach into
application payload, and only when it is readable. Where the monitor is
placed to capture these headers, and how sampling there affects detection,
is the subject of \cref{sec:data-acquisition}. What happens when payload is
encrypted is discussed in \cref{sec:feature-drift-encrypted-traffic}, and it
is one of the boundaries already fixed in \cref{sec:scope-limitations}.

% ================================================================
% 2.2 Data acquisition
% ================================================================
\section{Data acquisition}
\label{sec:data-acquisition}

Headers must be captured at an Observation Point before they can be
aggregated into packets, flows, connections or host statistics. This section
covers how a monitor is attached to the network, which standards export the
traffic as flow records, and how sampling at the point of capture trades
completeness for lower load.

\subsection{Packet capture, SPAN ports and TAPs}
\label{sec:packet-capture-mechanisms}

\emph{Full packet capture} records every packet crossing an Observation
Point, typically in the pcap format. It is the prerequisite for packet-level
analysis, since a signature-matching sensor needs the actual bytes and not a
derived record. How the sensor is attached determines what it can do with
them. \citet{scarfone2007nist80094} distinguish two placements. An
\emph{inline} sensor is deployed so that the traffic it monitors ``must pass
through it''. Its main motivation is to stop attacks by blocking traffic. A
\emph{passive} sensor ``monitors a copy of the actual network traffic''.
Passive techniques typically provide no reliable way to stop traffic from
reaching its destination. This capability trade-off is fixed by the
placement alone, whatever detection logic runs downstream. The same guide
lists three ways to obtain the copy. A spanning port, often called SPAN, is
``a port that can see all network traffic going through the switch''. A
network tap is ``a direct connection between a sensor and the physical
network media itself''. An IDS load balancer is ``a device that aggregates
and directs network traffic to monitoring systems''
\citep{scarfone2007nist80094}. The guide warns that a spanning port may not
see all traffic when the switch is under heavy load or badly configured, so
the copy is not always complete. Placement does not change what encrypted
payload reveals. A copy taken at a SPAN port or a tap carries the same
ciphertext as the original link.

\subsection{Flow export: NetFlow, IPFIX and sFlow}
\label{sec:flow-export-standards}

Once packets are captured, turning them into flow records requires an export
format on which the Metering Process and the collector agree. Cisco's
NetFlow introduced this model. Its early versions, including v5, export a
\emph{fixed} record format. A NetFlow v5 record has a fixed size of 48 bytes
\citep{sperotto2010flowbased}. NetFlow v9 replaced this with a
\emph{template}-based design. A Flow Record is no longer self-describing but
is interpreted against a separately exported Template, ``a collection of
fields, with corresponding descriptions of structure and semantics'', so new
field types can be added without changing the export format
\citep{claise2004rfc3954}. \gls{ipfix}, standardised by the \gls{ietf} in
RFC~7011, adopts the same template mechanism as an open, vendor-neutral
standard. This is why the flow definition used in this thesis comes from
IPFIX and not from NetFlow \citep{claise2013rfc7011}.

sFlow takes a different route. Instead of metering and exporting complete
flow records, an sFlow Agent combines statistical \emph{packet sampling}
with periodic polling of interface counters, and streams both, unaggregated,
to a central collector over UDP. Sampling works with a counter that
decrements with every packet, triggers a sample at zero and restarts from a
random value, so that the ratio of sampled to total packets converges to a
configured rate \citep{phaal2001rfc3176}. Unlike NetFlow and IPFIX, it does
not build flow records on the device. The collector receives samples and
counters, so its view of the traffic is statistical.

\subsection{Sampling and its effect on detection}
\label{sec:sampling-detection-effect}

Sampling is applied for the same reason as flow export, to keep monitoring
tractable on fast links. Two forms are common. \emph{Packet sampling}, on
which sFlow relies, selects a subset of packets \emph{before} flow records
are built, so under-sampled flows are formed from an incomplete view of the
traffic. \emph{Flow sampling} selects among flows, with each flow retained
with a random probability. \citet{sperotto2010flowbased} note that
sampling lowers the demands on the flow exporter but also makes the
detection of intrusions harder.

Sampling is not neutral for detection. \citet{brauckhoff2006samplingimpact}
took flow records from the backbone of a national network that contained
the Blaster worm outbreak, reconstructed a packet trace from them and
simulated packet sampling at rates from 1 in 10 to 1 in 1000. Byte and
packet counts are not disturbed, because they can be recovered by scaling
with the inverse of the sampling rate. Flow counts are grossly biased. A
scanning worm creates many small flows, and small flows are likely to be
missed, so in flow counts the worm is ``virtually undetectable at higher
sampling rates''. Entropy metrics, which summarize the distribution of
header fields such as addresses, are affected much less. The worm stays
exposed in flow entropy even at 1 in 1000.

For this thesis the consequence is that the observation level alone does not
guarantee that an attack is visible. Whether it survives the sampling regime
in place at the Observation Point is a separate question, which this thesis
does not study: it reads the records of the benchmark datasets as published
(\cref{ch:implementation}) and samples no packets itself.

% ================================================================
% 2.3 Feature engineering for traffic
% ================================================================
\section{Feature engineering for traffic}
\label{sec:feature-engineering-traffic}

None of the observation levels can be fed directly to a detector, symbolic or
learned. A raw packet, flow or connection record must first become a
fixed-length numerical representation. This section covers four recurring
families of such representations: statistics of a single flow, aggregates
over a time or connection window, features that require inspecting payload,
and the effect of widespread encryption on all of them.

\subsection{Flow statistics}
\label{sec:flow-statistics-cicflowmeter}

The most direct family works on a single flow. \citet{sharafaldin2018cicids2017}
describe CICFlowMeter, a flow-based feature extractor that computes more than
80 statistical features from a pcap file. Flows are identified by the
5-tuple of source and destination address, source and destination port, and
protocol, the same Flow Key as before. The features include the duration of
the flow, the minimum, mean, maximum and standard deviation of the
inter-arrival time, the number of bytes and packets sent in each direction,
packet-length statistics, counts of TCP flags such as SYN, ACK and PSH, the
initial window size in bytes, and packet rates. The dataset produced with
it, CICIDS2017, is a widely used benchmark for flow-based detection. Such a
feature set is a clear example of the payload-free, flow-level
representation to which the detectors of this thesis are restricted on the
flow datasets, as fixed in \cref{sec:scope-limitations}.

\subsection{Time-window and connection-window aggregates}
\label{sec:time-window-connection-aggregates}

A second family looks at the connections around the current one.
\citet{tavallaee2009kddcup99} define \emph{traffic features} as statistics
over the connections that share the destination host or the service of the
current connection, computed on a sliding window. The \emph{time-based}
variant uses the last two seconds and yields, for instance, the number of
connections to the same host or service. The \emph{connection-based} variant
uses a window of 100 connections. Both are second-order features. They
assume that connection records already exist and compute statistics across
them, and they say nothing about a connection in isolation.

\subsection{Content and host-based features}
\label{sec:content-host-based-features}

A third family departs from header statistics altogether.
\citet{tavallaee2009kddcup99} define \emph{content features} as indicators
obtained by inspecting the data portion of a connection for signs that
domain knowledge associates with an attack, for example the number of
failed login attempts. The reason is that R2L and U2R attacks, unlike DoS
and probing, do not involve many connections in a short period. They are
embedded in the data portions of the packets and normally involve a single
connection. Unlike the two families above, content features can be computed
only when the payload is readable. They depend on the same packet-level
visibility discussed in \cref{sec:tcpip-stack-monitor-view}, and encryption
removes it.

\subsection{Feature drift and the encrypted-traffic problem}
\label{sec:feature-drift-encrypted-traffic}

Content features were defined on traffic from an era when payload was
routinely readable. On today's largely encrypted traffic that assumption
fails. \citet{moriarty2018rfc8404} document the effect of pervasive
encryption on network operators. Payload examination ``in the clear'' offers
filtering opportunities that encryption removes, and estimating the type of
traffic becomes ``more complex and less accurate''. With transport-layer
encryption the 5-tuple stays visible, so NetFlow-type data can still be
collected. Only tunnelling, such as IPsec in tunnel mode, restricts the
observable fields further, to a 2-tuple of the tunnel endpoints. What
survives is, in other words, the flow- and header-level information of the
first two families.

For feature engineering this has a direct consequence. A detector designed
and validated on older datasets rich in content features can silently lose
the features it depends on when deployed on encrypted traffic. This
distribution mismatch is why the thesis keeps content features apart, as
fixed in \cref{sec:scope-limitations}. They are used only on NSL-KDD, where
they come with the records. On the flow datasets and in the experiments
across datasets the feature space is flow- and header-level only. A result
that rests on content features is therefore read as specific to NSL-KDD. The
separation is a design choice and not a failure discovered at deployment. For a survey of how analysis
techniques adapt to encrypted traffic, see \citet{papadogiannaki2021encryptedsurvey}.

% ================================================================
% 2.4 Classical approaches to traffic analysis and intrusion detection
% ================================================================
\section{Classical approaches to intrusion detection}
\label{sec:classical-approaches}

Before learning-based detectors, network monitoring relied on a small set of
classical methods. NIST distinguishes three detection methodologies:
signature-based, anomaly-based and stateful protocol analysis
\citep{scarfone2007nist80094}. This section follows that classification and
adds threshold and rate-limiting defences, which act on traffic rates and not
on attack patterns. Each family makes a different assumption about what an
attack looks like, and the differences show what later chapters try to
improve.

\subsection{Signature-based detection}
\label{sec:signature-rulesets}

Signature-based detection was defined in \cref{sec:signature-not-enough}.
Here the question is how the rules are built. \citet{roesch1999snort}
presents Snort as a libpcap-based packet sniffer and logger that can serve as
a lightweight \gls{nids}, with rule-based logging that performs content
pattern matching. A Snort rule has a header and options. The header gives the
protocol, the direction, and the addresses and ports. The options, written in
parentheses, add further tests, and all of them must be true for the rule to
match. The most basic rules contain only protocol, direction and port. A
matching rule triggers one of three actions: pass, log or alert. Snort itself
has three subsystems: a packet decoder, a detection engine and a logging and
alerting subsystem. New rules are written by running an exploit on a test
network, recording the traffic, finding a unique signature in it and
condensing that signature into a rule. The original paper also lists
limitations of that first version, among them the lack of TCP stream
reassembly.

Suricata rules have the same overall shape: an action, a header with
protocol, addresses, ports and direction, and options that define the
specifics of the rule \citep{suricata2026rules}. Zeek is built differently.
It was first published as Bro \citep{paxson1998bro} and renamed Zeek in 2018
\citep{zeek2018rename}. In Bro an event engine reduces the stream of packets
to higher-level network events, and an interpreter runs a specialized
scripting language in which a site's security policy is written as event
handlers \citep{paxson1998bro}. Detection is then expressed as scripts over
events and not as a list of patterns. For this thesis the important point is
that rules with content options need payload, while the most basic rules use
only header fields.

\subsection{Specification and protocol-anomaly detection}
\label{sec:specification-protocol-anomaly}

NIST calls this family stateful protocol analysis. It compares predetermined
profiles of generally accepted benign protocol activity, for each protocol
state, against observed events \citep{scarfone2007nist80094}. The profiles are
not learned on the local network. They are universal profiles developed by
vendors, which specify how a protocol should and should not be used. A
signature describes an attack. A specification describes normal protocol
behavior, so an attack is found when it breaks the protocol, even if no
signature for it exists. The analysis follows the state of each connection,
which TCP builds at the transport layer, as described in
\cref{sec:tcpip-stack-monitor-view}. NIST lists three limitations. The
analysis is resource-intensive because of its complexity. It cannot detect
attacks that stay within accepted protocol behavior. And the protocol model
used by the detector may conflict with the way the protocol is implemented
\citep{scarfone2007nist80094}. A specification is a set of explicit rules
written by experts, which makes it close in spirit to the domain knowledge
that later chapters encode as logical axioms.

\subsection{Statistical anomaly detection}
\label{sec:statistical-anomaly-detection}

NIST defines anomaly-based detection as the comparison of a profile of normal
behavior with current activity. The profiles describe users, hosts, network
connections or applications, are built by monitoring typical activity over a
period of time, and are compared with the current activity through
statistical methods and thresholds. NIST also lists the weaknesses. A profile
can inadvertently include malicious activity. Accurate profiles are hard to
build. And benign activity that deviates strongly from the profile often
produces many false positives \citep{scarfone2007nist80094}. Four
statistical techniques are common in the traffic-analysis literature.

\emph{Entropy} summarizes how concentrated or dispersed the distribution of a
packet feature is. \citet{lakhina2005mining} apply it to source and
destination addresses and ports in flow traces. Scans and small denial of
service attacks may barely change the traffic volume but change these
distributions. In the anomalies the authors found, the port scans, network
scans and point-to-multipoint transfers were not detected with volume
metrics and were found only with entropy. The robustness of entropy to
sampling was seen in \cref{sec:sampling-detection-effect}.

\emph{PCA} separates regular from irregular traffic. \citet{lakhina2004diagnosing}
apply it to traffic counts on the links of a network, in ten-minute bins. The
first principal components span a normal subspace that captures regular
patterns such as the daily cycle. The rest is a residual subspace, and a
volume anomaly shows up as a large residual. The authors note that the method
finds anomalies that are significant at the ten-minute scale, and that
anomalies are harder to detect in flows with large variance. Their later
work extends the method to several features at once \citep{lakhina2005mining}.

\emph{CUSUM} is a cumulative-sum test for a persistent change in a series.
\citet{wang2002detecting} use a non-parametric version to detect SYN
flooding. They monitor the difference between the number of SYN packets and
of FIN and RST packets, which is small in normal operation and grows when
spoofed SYN packets flood the network. The detector needs no per-connection
state, only counters for the three packet types. The authors point out that
an attacker who knows the scheme can defeat it by flooding a mixture of SYN
and FIN or RST packets.

\emph{EWMA} is an exponentially weighted moving average. \citet{muenz2008control}
apply Shewhart, CUSUM and EWMA control charts, not to the counts themselves
but to the prediction errors of exponential smoothing and Holt-Winters
forecasts. The data are byte, packet and flow counts from unsampled NetFlow
v5 records of an ISP backbone, aggregated in intervals of 300 seconds. An
alarm is raised when the statistic leaves limits defined as multiples of the
estimated standard deviation. Among the alarms the authors found a mailbomb,
a password guessing attempt and a network failure, but most of them were
caused by legitimate traffic, mainly large downloads from web or file
servers.

The four methods share one design: a statistic computed on the traffic, a
model of normal values and a threshold. None of them needs labelled attack
examples. They also share the weakness NIST names, since any strong benign
deviation raises an alert. This is the base-rate problem taken up in
\cref{sec:evaluation-practice}.

\subsection{Threshold and rate-limiting defences}
\label{sec:threshold-rate-limiting}

A threshold rule raises an alert, or drops traffic, when a count or a rate
exceeds a fixed limit, for example the number of packets per second from one
source. It is the statistical idea of the previous subsection with the limit
set by an operator instead of being estimated from data. This is the
weakness of the approach. \citet{handley2006rfc4732} state that in principle
it is not possible to distinguish a sufficiently subtle denial of service
attack from a flash crowd, that is, a legitimate surge of users. A limit low
enough to stop the attack also hits legitimate traffic, and a limit high
enough to spare legitimate users lets a subtle attack through. The same RFC
notes that filtering at the victim server does not help against attacks that
saturate a link, because the traffic has already crossed the scarce link, so
filtering has to happen closer to the source. Attacks designed to stay under
the limits make the problem worse. StormWall reports that reconnaissance
attacks, which usually last up to 15 minutes at relatively low intensity,
often stay below filtering thresholds \citep{stormwall2026forecast}. A
fixed numerical limit is also the simplest example of an explicit domain
rule given to a detector, which the neurosymbolic approach of this thesis
generalizes.

% ================================================================
% 2.5 Standards, taxonomies and reference architectures
% ================================================================
\section{Standards, taxonomies and reference architectures}
\label{sec:standards-taxonomies}

The methods of the previous section leave three questions open. How does a
detector fit into a deployment, how are the attacks it reports named, and
how do its alerts reach other tools? Standards and taxonomies answer them.
They detect nothing themselves. This section covers the ones a reader meets
next to the detectors of this thesis.

\subsection{NIST SP 800-94 and ISO/IEC 27033}
\label{sec:nist-iso-architectures}

NIST SP 800-94 is the source of the detection methodologies used so far. It
is also an architecture guide. Its abstract names four classes of intrusion
detection and prevention systems, or IDPS: network-based, wireless, network
behavior analysis and host-based. A typical deployment has four components:
sensors or agents, a management server, a database server and a console.
The guide lists four capabilities: information gathering, logging, detection
and prevention \citep{scarfone2007nist80094}. The class closest to this
thesis is network behavior analysis. The guide describes it as examining
network traffic for threats that produce unusual traffic flows, and names
\gls{ddos} attacks as an example \citep{scarfone2007nist80094}. The guide
dates from 2007. NIST drafted a revision in 2012 but never finalised it and
has retired the draft \citep{scarfone2007nist80094}.

ISO/IEC 27033 is a series on network security. Its first part gives an
overview of network security and related definitions and is a road map to
the other parts. The catalogue description speaks of risks, controls and
network security architecture \citep{iso2015iec27033}. The thesis does not
use it further.

\subsection{MITRE ATT\&CK as a vocabulary of behaviour}
\label{sec:mitre-attack-vocabulary}

\citet{strom2020attack} describe ATT\&CK as a knowledge base of adversary
tactics and techniques based on real-world observations. A tactic is the
tactical objective of the adversary, the why. A technique is how the
adversary reaches that objective by performing an action, the how. A
procedure is a specific implementation of a technique used by adversaries.
The first model was created in 2013 and focused mainly on Windows
enterprise environments. The ATT\&CK website now lists three matrices,
Enterprise, Mobile and ICS, and its current version is 19.2, released in
April 2026 \citep{mitre2026attackversions}.

Network flooding appears as technique T1498, Network Denial of Service,
under the tactic Impact. Its sub-techniques are Direct Network Flood and
Reflection Amplification \citep{mitre2025t1498}. One of the detection
strategies attached to it watches for executables or scripts that generate
a large volume of outbound traffic \citep{mitre2025t1498}. That is a view
from the host that sends the traffic and not from the link that receives it.
ATT\&CK therefore gives the attack families of \cref{ch:attacks-ddos} a name
in a shared vocabulary, but it gives no flow-level conditions that a network
detector could test. In this thesis it is a label space and not a rule set.

\subsection{Alert formats and exchange}
\label{sec:alert-formats}

Once a detector raises an alert, it has to be passed on. NIST lists e-mail,
pages, messages on the IDPS user interface, SNMP traps, syslog messages and
user-defined programs or scripts as delivery methods
\citep{scarfone2007nist80094}. Four specifications cover what is delivered
and how it travels.

IDMEF defines a data model for information exported by intrusion detection
systems, with an XML implementation. It has two message types, Alert and
Heartbeat. An Alert carries, among other things, the analyzer, the
classification, the source and the target. The RFC has Experimental status
and is not a candidate for an Internet Standard \citep{debar2007rfc4765}.

STIX 2.1 is an OASIS Standard, approved in June 2021. It is a language for
sharing cyber threat intelligence in machine-readable form. Information about
attacks, attackers and vulnerabilities is represented as linked objects such
as malware, threat actors and attack patterns \citep{jordan2021stix}.
TAXII 2.1 carries this content. It is an application layer protocol that runs
over HTTPS, with Collections, where consumers request data from a producer,
and Channels, for publish and subscribe. The specification reserves Channels
for future versions \citep{jordan2021taxii}.

CEF is the event format of ArcSight, described in a document of its
vendor and not in a standard of an independent body. A message has seven
header fields separated by pipes, namely version, device vendor, device
product, device version, event class ID, name and severity, followed by an
extension of key-value pairs. It uses syslog as transport
\citep{opentext2023cef}.

The four differ in scope. IDMEF describes an alert from one detector. STIX
describes threat knowledge and TAXII moves it. CEF is a compact event record.
Comparing detectors needs none of them. They matter for what happens to the
output afterwards. ATT\&CK and threat-intelligence feeds matter earlier, as a
possible source of formal knowledge for the detector itself, which is listed
as future work in \cref{ch:conclusions}.

% ================================================================
% 2.6 Evaluation practice in NIDS and its recurring flaws
% ================================================================
\section{Recurring flaws in NIDS evaluation}
\label{sec:evaluation-practice}

Every method in this chapter is judged on labelled traffic, and the way that
judgement is made has its own failure modes. They are old, well documented
and still common. The benchmark of this thesis has to avoid each of them, so
this section states them. It builds on the two critiques of
\cref{ch:introduction}, \citet{sommer2010outside} and
\citet{arp2022dosdonts}, and looks at what they mean for measurement.

\subsection{Base-rate fallacy and the cost of false positives}
\label{sec:base-rate-false-positives}

Attacks are rare in most traffic. A detector works on data with a strong
class imbalance, and its false positive rate cannot be read alone.
\citet{axelsson2000baserate} applied Bayes' theorem to intrusion detection.
For an operator the useful quantity is the probability that an alarm marks a
real intrusion:
\[
P(I \mid A) = \frac{P(I)\,P(A \mid I)}
  {P(I)\,P(A \mid I) + P(\neg I)\,P(A \mid \neg I)},
\]
where $I$ is an intrusion and $A$ an alarm. When $P(I)$ is small, the second
term of the denominator dominates unless $P(A \mid \neg I)$ is very small
too. In Axelsson's example an intrusion accounts for about $2 \cdot 10^{-5}$
of the events. Even with a detection rate of 1.0, the false alarm rate has to
be about $10^{-5}$ for $P(I \mid A)$ to reach 66\%. He concludes that the
false alarm rate limits the performance of an intrusion detection system
\citep{axelsson2000baserate}.

\citet{arp2022dosdonts} make the same point for machine learning. With a
class ratio of 1:100, a detector with 99\% true positives and 1\% false
positives produces 100 false positives for every 99 true positives, so about
half of its alarms are false. A rarer attack class makes it worse. Take one
attack flow in 10{,}000 and the same rates. In one million flows there are
100 attack flows, of which 99 are detected, and 9{,}999 benign flows are
flagged as attacks. About 1\% of the alarms are real. These figures are
illustrative and the calculation is ours. The base rate is also not fixed.
During a flood, attack flows can outnumber benign ones, and the same
detector then looks far better. A precision figure has meaning only together
with the prevalence of the data it was computed on.

Two costs follow. The first is analyst time. \citet{sommer2010outside} list
the high cost of errors among the difficulties of applying machine learning
to intrusion detection. A false positive costs the time of an analyst who
examines the reported incident, and, citing Axelsson, even a very small rate
of them can make a detector unusable. A false negative has its own cost,
since a single compromised system can undermine the IT infrastructure. The
second cost falls on users. An inline sensor can block traffic
\citep{scarfone2007nist80094}, and every false positive then denies service
to a legitimate user.

Accuracy hides both costs. \citet{arp2022dosdonts} show that on data with a
1:100 ratio a ROC curve can look excellent while the precision-recall curve
shows poor performance. They recommend precision, recall and precision-recall
curves for the detection of rare events. A benchmark that reports accuracy
alone would repeat the flaw. The metric set of the benchmark is fixed in
\cref{ch:benchmarking-methodology}.

\subsection{Imbalance, leakage and in-distribution bias}
\label{sec:imbalance-leakage-bias}

Three flaws concern the data and the way it is split. They are separate, and
each one inflates the result.

\emph{Class imbalance} is the property behind the previous subsection. A
dataset is a fixed sample, and \citet{arp2022dosdonts} note that researchers
usually do not know the true distribution of the data, so a fixed dataset may
differ from real conditions. The class ratio is one part of that difference.
The benchmark therefore treats the degree of imbalance as a variable, and
stronger imbalance is one of its stress scenarios.

\emph{Temporal leakage} lets information from the period under test reach
training. \citet{arp2022dosdonts} call it temporal snooping and define it as
ignoring the time dependencies within the data. In security the distribution
is not stationary, because attacks evolve. They recommend that dataset splits
take these dependencies into account and that test data be kept apart until
the final evaluation. A random split of one capture into training and test
sets ignores time, since it mixes earlier and later traffic.

\emph{In-distribution bias} is the term used here for a score that holds
only on the collection where it was measured. Training and test data come
from the same source, so the score measures fit to that source and not
detection. The closest pitfalls in \citet{arp2022dosdonts} are sampling bias
and lab-only evaluation. The KDD Cup 99 data are an example. The NSL-KDD
page reports that removing duplicate records cut the training set by 78.05\%
and the test set by 75.15\%. It explains that the redundancy biases learners
towards frequent records and the evaluation towards methods that do well on
them \citep{cic2009nslkdd}. NSL-KDD removes the duplicates. It does not
change the age of the traffic. The KDD Cup data derive from the 1998 DARPA data, which
\citet{sommer2010outside} already judged too old and too synthetic for
current studies. A second example is the case study of
\citet{arp2022dosdonts} on the Kitsune detector. On the Mirai dataset a
simple boxplot baseline outperformed it at low false positive rates. The
authors trace this to attack signals in the dataset that are easy to see and
not representative of real attacks.

The remedies point the same way. \citet{sommer2010outside} ask for
evaluation on data from several networks. \citet{arp2022dosdonts} suggest
building several estimates of the true distribution, analysing each one
separately and not mixing data from incompatible sources. Cross-dataset
transfer, one of the stress scenarios of the benchmark, follows from this.

\subsection{Two critiques of learning-based detection}
\label{sec:closed-world-dos-donts}

The two papers are twelve years apart and ask different questions.
\citet{sommer2010outside} explain why intrusion detection is a hard task for
machine learning. They name the closed world assumption, the cost of errors,
the semantic gap, the diversity of traffic and the difficulty of evaluation.
Their recommendations are to understand the threat model, keep the scope
narrow, justify the choice of features, reduce the cost of errors and
evaluate on real traffic from several networks with independent ground truth.
\citet{arp2022dosdonts} look at how studies go wrong. They define ten
pitfalls, from sampling bias to an inappropriate threat model, and give
recommendations for each.

They overlap in three places. Both ask for an explicit threat model. Both
raise the cost of false positives, as the cost of errors in one and the base
rate fallacy in the other. Both ask for evaluation data closer to real
conditions, as several networks in one and the pitfalls of sampling bias and
lab-only evaluation in the other. This overlap is our reading. It gives a
checklist for the benchmark. The design of the comparison takes it as a
constraint: matched baselines, several datasets, cross-dataset transfer,
imbalance as a controlled variable and metrics that stay informative under
imbalance.

The critiques have limits. \citet{sommer2010outside} argue mostly at an
intuitive level and do not prove their points formally. They list the
semantic gap and the diversity of traffic among the difficulties. This thesis
asks whether explicit domain knowledge helps with them. It does not assume the
answer, and the checklist applies to its own experiments as much as to the
literature it compares.

% ================================================================
% 2.7 Summary
% ================================================================
\section{Summary}
\label{sec:ch2-summary}

This chapter fixed what a detector in this thesis sees and how its results
are judged. The level of observation and the capture point decide which
attacks are visible. Packets, flows, sessions and host aggregates offer
different units of analysis, and sampling can hide an attack even at the
right level. Feature families range from flow statistics to content features,
and encryption removes the content features. On the flow datasets the
detectors compared here are therefore limited to flow- and header-level features, and content features are used only on
NSL-KDD.

The classical methods each rest on something explicit. A signature describes
an attack, a specification describes normal protocol behavior, a statistic
models normal traffic and a threshold sets a limit. Each fails in a known
way. Signatures miss unknown attacks, statistics and thresholds raise alarms
on benign deviations, and a subtle attack cannot in principle be told from a
flash crowd. NIST SP 800-94 gives the reference classification, ATT\&CK a
vocabulary for naming attacks, and IDMEF, STIX, TAXII and CEF formats for
passing information on.

Evaluation has its own recurring flaws: the base-rate fallacy, class
imbalance, temporal leakage and in-distribution bias. Two things carry
forward. Fixed rules and limits are the simplest explicit domain knowledge,
and neurosymbolic detection generalizes them. The checklist drawn from
Sommer and Paxson and from Arp et al. applies to the benchmark of this
thesis.
